\section{选题背景与意义}

\subsection{选题背景}

毛泽东在莫斯科大学会见中国留学生和实习生时指出：“世界是你们的，也是我们的，但是归根结底是你们的。你们青年人朝气蓬勃，正在兴旺时期，好像早晨八九点钟的太阳。希望寄托在你们身上。”\footcite{Mao1992Youth}这一论断深刻揭示了青年群体在推动社会发展中的重要地位。作为青年群体中的中坚力量，大学生的成长成才直接关系到国家现代化建设与民族复兴的进程。从个体成长规律来看，大学生正处于价值观确立与定型的关键期\footcite{CPCC2017_Plan}，其价值取向既关乎个人发展，也关系到未来整个社会的精神面貌\footcite{Xi2014}。引导大学生扣好人生的“第一粒扣子”、自觉树立和践行社会主义核心价值观，构成高校落实立德树人根本任务、培养德智体美劳全面发展的社会主义建设者和接班人的关键所在\footcite{CPCC2017_Ipe,Xi2018_Speech}。围绕高校育人与青年价值观塑造的战略需求，近年来党和国家作出一系列系统部署：2021年，《关于新时代加强和改进思想政治工作的意见》明确要求把思想政治工作贯穿于教育教学与日常管理全过程，贴近青年实际做好引导与关怀\footcite{CPCC2021_Yijian}；2024年，党的二十届三中全会对完善立德树人机制、健全培育和践行社会主义核心价值观制度机制作出系统安排\footcite{CPCC2024_Jueding}；2025年印发的《教育强国建设规划纲要（2024—2035年）》进一步提出坚持显性教育与隐性教育相统一，拓展实践育人与网络育人空间\footcite{CPCC2025_Gangyao}。落实上述部署，必须深入把握大学生价值观形成的现实机制。马克思在《关于费尔巴哈的提纲》中指出：“人的本质不是单个人所固有的抽象物，在其现实性上，它是一切社会关系的总和。”\footcite{MarxEngels2009}这一经典论述表明，大学生对马克思主义价值观和社会主义核心价值观的理解与认同深受现实社会关系与日常人际交往经验的塑造，亦为本研究确立了基本考察视角，即从大学生所处的具体社会关系出发，剖析日常交往方式变迁如何作用于其价值观。


在数字时代，随着经济社会发展、社会分工细化以及移动互联网的普及，当代大学生的日常交往方式正在发生显著变化。社交平台与即时通讯工具大幅降低了跨专业、跨高校乃至跨地域建立联系的门槛，大学生的日常社交在线上线下交织的多元场域中广泛展开，呈现出按需匹配、流动聚散的新样态。在这一交往图景中，大量人际互动围绕共同兴趣、具体任务或阶段性需求形成，普遍具有情感投入较浅、双方对彼此背景了解有限、承担的长期责任较轻等特征，在社会网络意义上构成了典型的弱连接交往，具体表现为自习搭子、临时项目伙伴、兴趣社群成员以及跨校求职互助同伴等。客观来看，弱连接交往借助跨群体的横向联结，能够帮助大学生获取异质性信息与多元生活经验，在拓宽认知视野、便利阶段性协作以及降低日常社交维系成本方面具有一定积极作用。

然而，当这种情感投入较浅、按需聚散的交往方式在大学生日常生活中不断扩张，并逐步挤占需要长期投入的关系空间时，其内在局限便会转化为现实的价值隐忧。一方面，由于弱连接交往对象之间缺乏全面了解，日常互动中集中呈现的片面生活片段、竞争结果与情绪化看法，容易使部分大学生在自我评价与消费休闲中滋生比较焦虑，在看待公共事务时趋于片面与偏激；另一方面，按即时需求建立联系、随任务结束而终止往来的低义务相处习惯，容易助长将专业学习窄化为短期功利筹码、将日常人际往来简化为得失交换的实用倾向，并在时间与精力上挤占深度学习钻研、自我修养以及维系稳定情感联结的投入。上述微观层面的观念偏移在日常生活中持续累积，与马克思主义关于人的全面发展、崇尚诚实劳动和集体互助的价值立场产生张力，亦同社会主义核心价值观所倡导的爱国、敬业、诚信、友善等育人目标形成现实偏差。推进新时代大学生思想政治教育，必须坚持因事而化、因时而进、因势而新\footcite{Xi2022}。面对大学生交往形态向跨边界、流动性演进的现实变化，亟须突破以往偏重固定群体关系的考察视角，深入把握弱连接交往作用于大学生日常价值观的内在过程。

基于上述现实背景，本研究立足立德树人根本任务与思想政治教育学科视角，以弱连接交往对大学生价值观的影响为研究主题，系统考察弱连接交往在个人生活、人际交往、学习劳动与择业以及社会公共四个维度引发的观念变化，深入剖析浅层互动规则传导至日常价值观的微观机理与深层成因，揭示其对马克思主义价值观及其当代表现——社会主义核心价值观培育带来的现实挑战，并结合理论研究与网络文本质性分析探寻具有针对性的应对策略。
\subsection{选题意义}

\subsubsection{理论意义}

第一，从日常交往方式变迁切入，有助于拓展大学生价值观影响因素的理论分析。马克思关于人的本质是一切社会关系总和的论述表明，大学生价值观的形成与演变根植于现实的人际交往实践。既有关于弱连接交往的研究多集中于信息传递、求职互助和情绪慰藉等实用功能，对这类浅层交往如何作用于大学生价值观探讨较少。本研究立足马克思主义社会关系理论，重点考察弱连接交往与大学生价值观之间的内在关联，剖析浅层互动规则引发价值观变化的具体过程，为从日常社会关系维度把握大学生价值观演变提供新的分析维度。

第二，聚焦弱连接交往引发的价值变化，有助于拓展大学生社会主义核心价值观培育的理论视角。社会主义核心价值观的培育植根于大学生的日常生活实践，需要及时回应人际交往形态的时代演进。既有关于大学生社会主义核心价值观培育的研究，多侧重于宏观理论阐释、思政课教学及常规教育载体，对数字时代按需聚散、情感投入较浅的弱连接交往场域关注有限。本研究立足当代大学生的具体交往场景，考察弱连接交往如何作用于大学生对核心价值要求的理解与践行，为推动社会主义核心价值观融入大学生日常交往提供贴近现实的理论参照。
\subsubsection{实践意义}

第一，通过梳理大学生在弱连接交往中呈现的真实心态与价值观困惑，为高校一线教师和辅导员把握学情、选取贴近生活的讨论议题提供经验参考。跨校备考互助、兴趣社群交流及校园“搭子”往来中滋生的攀比焦虑与关系疏离，仅凭常规日常管理较难及时察觉。本研究借助对网络社区交往叙述与互动回复的系统分析，具体呈现大学生在流动社交中的现实关切与价值观波动，能够为一线思政课教师补充贴近生活的课堂案例，也为辅导员开展有针对性的日常沟通提供鲜活素材。

第二，针对弱连接交往中的具体矛盾提出可操作的引导思路，为高校基层改善校园活动设计以及大学生理性调适日常交往提供现实参考。大学生在参与跨班级、跨高校弱连接交往时，如何平衡个人边界诉求与集体协作责任，是课堂讨论、社团运作与社会实践共同面对的实际问题。本研究结合课程教学、校园协作项目与校外实践场景，探寻在尊重个人边界的前提下将主流价值融入日常互动的具体做法\footcite{Shen2021_narrative}，既有助于一线组织者改进集体活动的参与体验，也能为大学生理性看待浅层社交、缓解功利焦虑与人际疏离提供切实借鉴。
